\subsection{Margin Call (2011)}

La película presenta el momento en que una institución financiera descubre que los modelos de riesgo utilizados para evaluar su portafolio ya no describen correctamente la realidad del mercado. A partir de ese momento, la historia se centra en las decisiones que deben tomarse cuando los cálculos indican que las pérdidas potenciales podrían superar el capital de la empresa. Más que narrar el colapso financiero en sí, la película muestra el instante en que los actores del sistema comprenden que existe un problema estructural y deben reaccionar rápidamente para evitar la quiebra.

Un recurso narrativo relevante es la subtrama de la muerte del perro, que introduce un tono emocional y melancólico en la historia. Este elemento puede interpretarse simbólicamente como una representación de la “muerte de la confianza” entre instituciones financieras durante la crisis, ya que , la confianza constituye uno de los fundamentos más importantes del sistema financiero. El valor del dinero fiduciario, de los activos financieros y de instrumentos como los bonos, préstamos o derivados depende en gran medida de la creencia compartida de que las promesas de pago serán cumplidas en el futuro. Cuando esta confianza se deteriora, las instituciones dejan de prestarse entre sí, los mercados pierden liquidez y el sistema comienza a paralizarse. 

La película también muestra la problemática sistémica que forma parte de las bases del funcionamiento de la economía. Las decisiones se toman a partir de variables cuantitativas y métricas monetarias: pérdidas potenciales, capital disponible, precios de mercado o impacto de liquidar posiciones. La única forma de medir el problema es el dinero; en los modelos y discusiones no aparecen conceptos como el valor moral o la confianza, sino magnitudes financieras que permiten tomar decisiones dentro de las reglas del mercado.

Desde esta perspectiva, las acciones que toma la institución no necesariamente pueden interpretarse como incorrectas. Ante una situación similar, es probable que la mayoría de los agentes del mercado actuaran de manera comparable para proteger la supervivencia de su organización. La venta acelerada de activos problemáticos, incluso si esto transfiere el riesgo a otros participantes del mercado, responde a una lógica estratégica que busca evitar la insolvencia de la empresa.

Lo que resulta más criticable no es tanto lo que se hace, sino la forma en que se hace. La película muestra una notable indiferencia institucional hacia el personal, especialmente en los despidos, que se realizan de manera abrupta y con poca consideración hacia las personas afectadas. Algo similar ocurre con la venta de los instrumentos financieros problemáticos, se supone que son los CDOs: la decisión de venderlos puede entenderse dentro de la lógica de supervivencia de la firma, pero el modo en que se ejecuta, sabiendo que los activos son altamente riesgosos y presionando para colocarlos en el mercado antes que los demás, revela tensiones éticas dentro del sistema.

La película no presenta simplemente una historia de culpables individuales, sino un retrato de cómo funcionan las instituciones financieras en situaciones extremas. Las decisiones aparecen condicionadas por incentivos, modelos cuantitativos y dinámicas organizacionales que llevan a los actores a actuar de manera racional desde su propia posición, incluso cuando esas decisiones pueden contribuir a amplificar los problemas del sistema financiero en su conjunto.
